% ============================================================================
% CODE QUALITY
% ============================================================================

\section{Code Quality Analysis}
\label{sec:code-quality}

This section presents the results of automated code quality analysis, including complexity metrics and duplicate code detection.

\subsection{Complexity Analysis}

Complexity is measured using fan-out (outgoing calls) and fan-in (incoming calls) metrics. Functions with high fan-out are often doing too much and should be decomposed.

\subsubsection{Severity Levels}

\begin{table}[h]
\centering
\begin{tabular}{lll}
\toprule
\textbf{Severity} & \textbf{Fan-Out Threshold} & \textbf{Description} \\
\midrule
\critical{Critical} & $>$ 50 & Function calls far too many others \\
\warning{High} & $>$ 30 & Significant complexity, refactoring recommended \\
Medium & $>$ 10 (threshold) & Moderate complexity, monitor \\
\success{Low} & $\leq$ 10 & Acceptable complexity \\
\bottomrule
\end{tabular}
\caption{Complexity severity levels}
\label{tab:severity}
\end{table}

\subsubsection{High-Complexity Functions}

The analysis identified \textbf{17 functions} exceeding the recommended complexity threshold:

\begin{longtable}{lrrll}
\toprule
\textbf{Function} & \textbf{Fan-Out} & \textbf{Fan-In} & \textbf{Severity} & \textbf{Location} \\
\midrule
\endhead
\func{extract\_symbols} & 48 & 4 & \warning{HIGH} & \file{solidity.rs:151} \\
\func{extract\_symbols} & 48 & 4 & \warning{HIGH} & \file{typescript.rs:310} \\
\func{discover\_files} & 46 & 2 & \warning{HIGH} & \file{walker.rs:67} \\
\func{parse\_response} & 45 & 1 & \warning{HIGH} & \file{openai.rs:108} \\
\func{extract\_symbols} & 42 & 3 & \warning{HIGH} & \file{rust.rs:200} \\
\func{watch\_and\_index} & 41 & 1 & \warning{HIGH} & \file{index/mod.rs:428} \\
\func{build\_signature} & 39 & 2 & \warning{HIGH} & \file{typescript.rs:658} \\
\func{extract\_call\_edges} & 37 & 4 & \warning{HIGH} & \file{parser/mod.rs:276} \\
\func{run\_embed\_watch} & 37 & 1 & \warning{HIGH} & \file{main.rs:191} \\
\func{run\_query} & 37 & 1 & \warning{HIGH} & \file{main.rs:536} \\
\func{extract\_symbols} & 35 & 3 & \warning{HIGH} & \file{python.rs:186} \\
\func{extract\_edges} & 33 & 2 & \warning{HIGH} & \file{typescript.rs:430} \\
\func{run\_index} & 32 & 1 & \warning{HIGH} & \file{main.rs:258} \\
\func{index\_file} & 31 & 3 & \warning{HIGH} & \file{index/mod.rs:156} \\
\func{run\_search} & 31 & 1 & \warning{HIGH} & \file{main.rs:388} \\
\func{extract\_edges} & 31 & 2 & \warning{HIGH} & \file{python.rs:318} \\
\func{build\_signature} & 31 & 3 & \warning{HIGH} & \file{rust.rs:402} \\
\bottomrule
\caption{High-complexity functions requiring attention}
\label{tab:complexity}
\end{longtable}

\subsubsection{Analysis of Complex Functions}

\paragraph{\func{extract\_symbols} (Multiple Parsers)}
The \func{extract\_symbols} function in each language parser handles too many responsibilities:

\begin{itemize}
    \item Tree-sitter query execution
    \item Symbol metadata extraction
    \item Signature building
    \item Docstring extraction
    \item Visibility determination
    \item Parent-child relationship tracking
\end{itemize}

\textbf{Recommendation:} Decompose into smaller, focused functions:

\begin{lstlisting}[style=rustcode, caption={Proposed refactoring for extract\_symbols}]
fn extract_function_symbols(tree: &Tree, source: &str) -> Vec<Symbol>;
fn extract_class_symbols(tree: &Tree, source: &str) -> Vec<Symbol>;
fn build_symbol_signature(node: &Node, source: &str) -> String;
fn extract_docstring(node: &Node, source: &str) -> Option<String>;
fn determine_visibility(node: &Node) -> Visibility;
\end{lstlisting}

\paragraph{\func{discover\_files} (\file{walker.rs})}
Complex file walking logic combining:

\begin{itemize}
    \item Path resolution and normalization
    \item Ignore pattern matching
    \item Binary file detection
    \item Symlink handling
\end{itemize}

\textbf{Recommendation:} Extract helper functions for each concern.

\paragraph{\func{run\_query} (\file{main.rs})}
Large command dispatch function handling all query subcommands.

\textbf{Recommendation:} Use a command pattern or dispatch table:

\begin{lstlisting}[style=rustcode, caption={Command dispatch pattern}]
type QueryHandler = fn(&Database, &QueryArgs) -> Result<()>;

fn get_query_handlers() -> HashMap<&'static str, QueryHandler> {
    let mut handlers = HashMap::new();
    handlers.insert("callers", handle_callers as QueryHandler);
    handlers.insert("deps", handle_deps as QueryHandler);
    handlers.insert("impact", handle_impact as QueryHandler);
    handlers
}
\end{lstlisting}

\subsection{Duplicate Code Detection}

Duplicate detection uses hash-based comparison with Jaccard similarity. The analysis identified \textbf{13 duplicate pairs} with similarity $\geq$ 80\%.

\begin{longtable}{rll}
\toprule
\textbf{Similarity} & \textbf{Location 1} & \textbf{Location 2} \\
\midrule
\endhead
90.9\% & \func{extract\_edges} & \func{extract\_edges} \\
       & \file{python.rs:318} & \file{typescript.rs:430} \\
\midrule
88.2\% & \func{stream\_start} & \func{stream\_start} \\
       & \file{formatter.rs:106} & \file{formatter.rs:142} \\
\midrule
85.0\% & \func{wrap} & \func{wrap} \\
       & \file{formatter.rs:99} & \file{formatter.rs:135} \\
\midrule
82.9\% & \func{test\_extract\_calls} & \func{test\_extract\_calls} \\
       & \file{python.rs:877} & \file{rust.rs:614} \\
\midrule
81.0\% & \func{get\_outgoing\_edges} & \func{get\_incoming\_edges} \\
       & \file{schema.rs:355} & \file{schema.rs:370} \\
\bottomrule
\caption{Detected duplicate code pairs}
\label{tab:duplicates}
\end{longtable}

\subsubsection{Duplicate Categories}

\paragraph{Formatter Implementations}
The \code{Formatter} trait implementations share significant code structure. The \func{stream\_start} and \func{wrap} methods are nearly identical across formatters.

\textbf{Recommendation:} Use default trait methods or a macro:

\begin{lstlisting}[style=rustcode, caption={Default trait implementation}]
trait Formatter {
    fn format_name(&self) -> &'static str;
    fn file_wrapper(&self) -> (&'static str, &'static str);
    
    // Default implementation using the above
    fn stream_start(&self, tree: Option<&str>) -> String {
        let (start, _) = self.file_wrapper();
        format!("{}{}", start, tree.unwrap_or(""))
    }
}
\end{lstlisting}

\paragraph{Edge Query Functions}
The \func{get\_outgoing\_edges} and \func{get\_incoming\_edges} functions in \file{schema.rs} differ only in the SQL WHERE clause.

\textbf{Recommendation:} Parameterize the direction:

\begin{lstlisting}[style=rustcode, caption={Unified edge query function}]
pub enum EdgeDirection {
    Outgoing,
    Incoming,
}

pub fn get_edges(&self, symbol_id: &str, direction: EdgeDirection) 
    -> Result<Vec<Edge>> 
{
    let column = match direction {
        EdgeDirection::Outgoing => "source_id",
        EdgeDirection::Incoming => "target_id",
    };
    // Single implementation with parameterized query
}
\end{lstlisting}

\paragraph{Test Utilities}
Parser tests share similar patterns for testing call extraction.

\textbf{Recommendation:} Create shared test fixtures:

\begin{lstlisting}[style=rustcode, caption={Shared test fixture}]
// In tests/common/mod.rs
pub fn assert_call_edges(
    source: &str, 
    expected_calls: &[(&str, &str)],
    parser: impl Parser
) {
    let edges = parser.extract_edges(source);
    for (caller, callee) in expected_calls {
        assert!(edges.iter().any(|e| 
            e.source_name == *caller && e.target_name == *callee
        ));
    }
}
\end{lstlisting}

\subsection{Summary}

\begin{table}[h]
\centering
\begin{tabular}{lr}
\toprule
\textbf{Quality Metric} & \textbf{Count} \\
\midrule
Total functions analyzed & 446 \\
High complexity functions & 17 (3.8\%) \\
Critical complexity functions & 0 \\
Duplicate code pairs & 13 \\
Average fan-out & 6.0 \\
Maximum fan-out & 48 \\
\bottomrule
\end{tabular}
\caption{Code quality summary}
\label{tab:quality-summary}
\end{table}

The codebase is generally well-structured, with complexity issues concentrated in the parser modules and main command handlers. Addressing the identified duplicates would improve maintainability and reduce the risk of divergent implementations.
